\documentclass{article}
\usepackage[cm]{fullpage}
\usepackage{newtxmath}
\usepackage{newtxtext}
\usepackage{graphicx}
\usepackage{hyperref}
\usepackage{placeins}

\begin{document}
\title{2785209M - James Middleton - Assessed Exercise - COMPSCI4011 Operating Systems (H)}
\date{}
\maketitle

\section*{Ownership statement}
\begin{flushleft}
    This is my own work, except for the code that was provided from moodle in the test harness, which we were permitted to use.
\end{flushleft}

\section*{Overview}
\begin{flushleft}
    From testing with the test harness, I believe the Solution functions properly within the expected guidelines.
    It utilizes the DiskDevice and FreeSectorDescriptorStore to produce a well optimized solution.
    From running it side by side with the provided demo, it seems to produce an almost identical output, albeit maybe slightly slower.\\
    The solution uses a separate write and read queue (Full Duplex) to optimize the efficiency of the code and speed of the process. \\
    There is no use of malloc, and all memory is allocated using a static voucher pool.\\
    The BoundedBuffer code from Moodle was used.\\
    There are shutdown sentinels within the internal worker threads to close the threads during shutdown.
\end{flushleft}

\section*{Build Instructions}
\begin{flushleft}
    To compile, I had to edit the makefile to disable pie.
    \\
    \begin{verbatim}
        cc -no-pie -o my_demo $(OBJECTS) -lpthread
    \end{verbatim}
\end{flushleft}


\begin{figure}
    \includegraphics[width=\textwidth]{DiskDriverSequenceDiagram.png}
    \caption{Design Sketch}
\end{figure}
\end{document}